% 仅含本问 bibitem，由主参谋放入全文 thebibliography 环境。
% q1-data 的作者、中文题名和 DOI 取自 WaterNet 正文的“数据的引用”，
% 不将页面标题区 10.1888 的拼写沿用为 DOI；实际阅读为元数据镜像。
\bibitem{q1-data}
马耀明，马伟强，谢志鹏，王宾宾，陈学龙，韩存博.
西藏吉隆高寒山区地气相互作用逐小时综合观测数据集（2021--2026）[DS/OL].
国家青藏高原科学数据中心，2026.
\url{https://doi.org/10.11888/Atmos.tpdc.303600}.
WaterNet 元数据镜像[2026-10-24]：
\url{https://gwfnet.net/Metadata/Index/T-2026-08-28-p1f3e5qIqTEyibfPAp3smBeA}.

% 只作收缩协方差思想出处：本文固定 lambda=0.5，敏感性0.1/0.9，
% 采用单位阵目标；不是 Ledoit--Wolf 自动估计收缩强度或其最优性复现。
\bibitem{q1-shrinkage}
Ledoit O, Wolf M.
A well-conditioned estimator for large-dimensional covariance matrices[J].
Journal of Multivariate Analysis, 2004, 88(2): 365--411.
\url{https://doi.org/10.1016/S0047-259X(03)00096-4}.

% 可作自然类比的概念出处；本文自族排除、年/过程配额与局部距离
% 是本研究登记设计，不能归为 Lorenz 原文的直接算法。
\bibitem{q1-analog}
Lorenz E N.
Atmospheric predictability as revealed by naturally occurring analogues[J].
Journal of the Atmospheric Sciences, 1969, 26(4): 636--646.
\url{https://doi.org/10.1175/1520-0469(1969)26<636:APARBN>2.0.CO;2}.
